\documentclass{article}
\usepackage[top=2cm,bottom=2cm,left=3cm,right=3cm]{geometry}
\usepackage{enumitem}
\usepackage{xcolor}
\usepackage{minted}
\usepackage[most]{tcolorbox}
\tcbuselibrary{minted,breakable}
\usepackage{fontspec}
\usepackage{polyglossia}
\setdefaultlanguage{ukrainian}
\setotherlanguages{english}
\usepackage{Alegreya}
\usepackage{FiraMono}
\newfontfamily\cyrillicfont{Alegreya}
\newfontfamily\cyrillicfontsf{Alegreya Sans}
\newfontfamily\cyrillicfonttt{DejaVu Sans Mono}
\usepackage{sciffi-python}
\usepackage{sciffi-memo}

\usemintedstyle{emacs}

\newtcolorbox{CodeBox}[1][]{
    breakable,
    enhanced,
    colback=gray!10,
    colframe=gray!70,
    fonttitle=\bfseries,
    title=#1,
    sharp corners,
    boxrule=0.8pt,
    top=2mm,
    bottom=2mm,
    left=2mm,
    right=2mm,
    fontupper=\footnotesize,
}

\newtcolorbox{ResultBox}[1][]{
    breakable,
    enhanced,
    colback=gray!5,
    colframe=gray!50,
    fonttitle=\bfseries,
    title=#1,
    sharp corners,
    boxrule=0.6pt,
    top=1mm,
    bottom=1mm,
    left=2mm,
    right=2mm,
    fontupper=\footnotesize,
}

\begin{document}
    \begin{titlepage}
        \begin{center}
            {\Large
                МІНІСТЕРСТВО ОСВІТИ ТА НАУКИ УКРАЇНИ

                НАЦІОНАЛЬНИЙ ТЕХНІЧНИЙ УНІВЕРСИТЕТ УКРАЇНИ
                «КИЇВСЬКИЙ ПОЛІТЕХНІЧНИЙ ІНСТИТУТ ІМЕНІ ІГОРЯ СІКОРСЬКОГО»

                ФАКУЛЬТЕТ ІНФОРМАТИКИ ТА ОБЧИСЛЮВАЛЬНОЇ ТЕХНІКИ

                КАФЕДРА ІНФОРМАЦІЙНИХ СИСТЕМ ТА ТЕХНОЛОГІЙ
            }


            \vspace{60mm}
            {\large
                \textbf{ЗВІТ}

                \vspace{5mm}

                \textbf{До лабораторної роботи 2}

                \vspace{5mm}

                з дисципліни «Технології Розробки Програмного Забезпечення»
            }

            На тему «Музичний Магазин»

            варіант \textnumero 14 (\textnumero 4)
        \end{center}

        \vfill
        \hfill
        \begin{minipage}[t]{0.35\textwidth}
            \textbf{Виконав:}
        \end{minipage}

        \vfill

        \begin{center}
            {\bf
                Київ

                КПІ Ім. Ігоря Сікорського

                2025
            }
        \end{center}
    \end{titlepage}
    \newpage

    Тема: запити до даних БД.

    Мета: вивчити команди формування простих запитів до бази даних (SELECT).

    Навички що будуть здобуті: вміння формувати прості запити до бази даних.

    \section*{Хід виконання:}
    \begin{sciffi}{python}
        from pathlib import Path
        from subprocess import run, PIPE
        from collections.abc import Sequence
        from contextlib import contextmanager
        
        @contextmanager
        def env(name: str, *args: str):
            print(f"\\begin{{{name}}}" + "".join(args))
            try:
                yield
            finally:
                print(f"\\end{{{name}}}")

        def sh(*args: str | Path) -> str:
            return run(
                tuple(map(str, args)),
                stdout=PIPE,
                stderr=PIPE,
                text=True,
            ).stdout

        scripts = Path("./scripts/")

        flow: Sequence[tuple[Path, Sequence[str]]] = [
            (script, ("--no-echo",))
            for script in sorted(
                scripts.glob("*.sql"),
                key=lambda p: int(p.stem),
            )
        ]

        for script, args in flow:
            source = script.read_text(encoding="utf-8")
            output = sh("tsqlexec", *args, script)

            print(f"\\subsection*{{Запит з файлу: {script.name}}}")
            with env("CodeBox", "[SQL Source]"):
                with env("minted", "[breaklines,linenos]", "{sql}"):
                    print(source, end="")
            print()

            if not output.strip():
                continue

            print("Result:")
            with env("ResultBox", "[Output]"):
                with env("minted", "[breaklines]", "{text}"):
                    print(output, end="")
            print("\\vspace{5mm}")
    \end{sciffi}

    \newpage
    \section*{Контрольні питання}
    \begin{enumerate}[label=\arabic*., leftmargin=*]
        \item \textbf{Які частини в команді SELECT є обов’язковими?}
        
        Обов'язковими частинами команди \texttt{SELECT} є власне ключове слово \texttt{SELECT}, що вказує, які стовпці потрібно повернути, та ключове слово \texttt{FROM}, що вказує на таблицю (або таблиці), з якої вибираються дані.
        \begin{CodeBox}[Мінімальний запит]
        \begin{minted}[breaklines]{sql}
SELECT column1, column2
FROM table_name;
        \end{minted}
        \end{CodeBox}
        
        \item \textbf{Яка команда використовується для виключення повторів із результату запита?}
        
        Для виключення дублікатів рядків з результату запиту використовується ключове слово \texttt{DISTINCT}. Воно застосовується одразу після \texttt{SELECT}.
        \begin{CodeBox}[Приклад з DISTINCT]
        \begin{minted}[breaklines]{sql}
SELECT DISTINCT city
FROM Customers;
        \end{minted}
        \end{CodeBox}
        
        \item \textbf{За допомогою якої команди створюються групи?}
        
        Групи створюються за допомогою речення \texttt{GROUP BY}. Воно групує рядки, що мають однакові значення у вказаних стовпцях, в один сумарний рядок. Зазвичай використовується з агрегатними функціями (\texttt{COUNT}, \texttt{SUM}, \texttt{AVG} тощо).
        \begin{CodeBox}[Приклад з GROUP BY]
        \begin{minted}[breaklines]{sql}
SELECT country, COUNT(CustomerID)
FROM Customers
GROUP BY country;
        \end{minted}
        \end{CodeBox}
        
        \item \textbf{В якій послідовності в команді SELECT повинні розміщуватися команди ORDER BY, GROUP BY, WHERE?}
        
        Логічна послідовність виконання та синтаксичний порядок цих речень є таким:
        \begin{enumerate}
            \item \texttt{FROM} -- вказує на джерело даних.
            \item \texttt{WHERE} -- фільтрує рядки до групування.
            \item \texttt{GROUP BY} -- групує відфільтровані рядки.
            \item \texttt{HAVING} -- фільтрує групи після групування.
            \item \texttt{SELECT} -- вибирає фінальні стовпці.
            \item \texttt{ORDER BY} -- сортує кінцевий результат.
        \end{enumerate}
        Отже, правильний синтаксичний порядок: \texttt{WHERE}, потім \texttt{GROUP BY}, і в кінці \texttt{ORDER BY}.
        
        \item \textbf{В чому різниця між командами WHERE та HAVING?}
        
        Основна різниця полягає в тому, на якому етапі обробки запиту вони застосовуються:
        \begin{itemize}
            \item \textbf{\texttt{WHERE}} фільтрує \textbf{окремі рядки} \textit{до того}, як вони будуть згруповані. Умови в \texttt{WHERE} не можуть містити агрегатних функцій.
            \item \textbf{\texttt{HAVING}} фільтрує \textbf{групи рядків} \textit{після того}, як вони були створені за допомогою \texttt{GROUP BY}. Речення \texttt{HAVING} використовується для застосування умов до результатів агрегатних функцій.
        \end{itemize}
        
        \item \textbf{Чи завжди запит, побудований з використанням команди HAVING може бути змінений на еквівалентний запит без використання цієї команди?}
        
        Ні, не завжди. Якщо умова фільтрації базується на результаті агрегатної функції (наприклад, \texttt{SUM(Price) > 1000} або \texttt{COUNT(*) > 5}), то таку умову неможливо перенести в речення \texttt{WHERE}. \texttt{WHERE} працює з даними на рівні окремих рядків, до агрегації, тому воно "не знає" результат \texttt{SUM} чи \texttt{COUNT}. У таких випадках використання \texttt{HAVING} є обов'язковим.
        
        \item \textbf{Які команди дозволяють обмежувати кількість строк що виводяться?}
        
        Для обмеження кількості рядків у результаті запиту використовуються різні команди залежно від діалекту SQL:
        \begin{itemize}
            \item \textbf{T-SQL (SQL Server):} \texttt{TOP n} або \texttt{TOP n PERCENT}
            \begin{CodeBox}[Приклад]
                \begin{minted}{sql}
SELECT TOP 10 * FROM Products;
                \end{minted}
            \end{CodeBox}
            \item \textbf{MySQL, PostgreSQL:} \texttt{LIMIT n}
            \item \textbf{Oracle, Standard SQL:} \texttt{FETCH FIRST n ROWS ONLY}
        \end{itemize}
        
        \item \textbf{В чому відмінність між простим запитом та складним?}
        
        \begin{itemize}
            \item \textbf{Простий запит} зазвичай звертається до однієї таблиці, не використовує підзапити, об'єднання (\texttt{JOIN}), оператори \texttt{UNION}, \texttt{INTERSECT}, \texttt{EXCEPT} або складні агрегації.
            \item \textbf{Складний запит} включає одну або декілька з таких конструкцій: об'єднання кількох таблиць (\texttt{JOIN}), підзапити (вкладені \texttt{SELECT}), операції над множинами, використання складних функцій або агрегацій з групуванням.
        \end{itemize}
        
        \item \textbf{За допомогою яких ключових слів в команді SELECT виконується перевірка входження результату обчислення виразу в задану множину?}
        
        Для перевірки входження значення в множину використовується оператор \texttt{IN}.
        \begin{CodeBox}[Приклад з IN]
        \begin{minted}[breaklines]{sql}
SELECT * FROM Customers
WHERE Country IN ('Germany', 'France', 'UK');
        \end{minted}
        \end{CodeBox}
        Також для цього можуть використовуватися оператори \texttt{ANY} та \texttt{ALL} у поєднанні з підзапитом.

        \item \textbf{Яка функція використовується для обрахунку кортежів, що відповідають заданій умові?}
        
        Для підрахунку кількості рядків (кортежів) використовується агрегатна функція \texttt{COUNT()}.
        \begin{itemize}
            \item \texttt{COUNT(*)}: підраховує всі рядки в групі.
            \item \texttt{COUNT(column\_name)}: підраховує рядки, де вказаний стовпець має значення, відмінне від \texttt{NULL}.
        \end{itemize}
        \begin{CodeBox}[Приклад з COUNT]
        \begin{minted}[breaklines]{sql}
SELECT COUNT(*) FROM Orders WHERE OrderDate >= '2023-01-01';
        \end{minted}
        \end{CodeBox}
        
        \item \textbf{Яка із команд команди SELECT накладає обмеження на відбір груп?}
        
        Обмеження на відбір груп, створених за допомогою \texttt{GROUP BY}, накладає речення \texttt{HAVING}.
        
        \item \textbf{Чи можливо використовуючи команду SELECT реалізувати операцію ділення?}
        
        Так, можливо реалізувати операцію реляційного ділення, хоча в SQL немає спеціального оператора для цього. Зазвичай це робиться за допомогою комбінації інших операторів, наприклад, через подвійне заперечення з \texttt{NOT EXISTS} або за допомогою \texttt{GROUP BY} та \texttt{COUNT}. Класичний приклад — "знайти студентів, які відвідують \textit{усі} курси з певного списку".
        
        \item \textbf{В яких випадках неможна для здійснення операції вибору використовувати команду WHERE?}
        
        Команду \texttt{WHERE} не можна використовувати, якщо умова фільтрації залежить від результату агрегатної функції (\texttt{SUM}, \texttt{AVG}, \texttt{MAX}, \texttt{MIN}, \texttt{COUNT}). \texttt{WHERE} виконується до обчислення агрегатів. У таких випадках необхідно використовувати \texttt{HAVING}.
        
    \end{enumerate}
\end{document}
